\chapter{Cardinals i aritmètica cardinal}\label{cap:cardinals}

\begin{objectius}
\apartat{Prerequisits} Del capítol~\ref{cap:cantor}: comparació i equipotència, Schröder--Bernstein, Cantor i la diagonal. Del capítol~\ref{cap:zorn}: Hartogs, el teorema del bon ordre i les equivalències de l'elecció. Del capítol~\ref{cap:families}: sumes i productes de famílies. Dels capítols~\ref{cap:valors}, \ref{cap:parts} i~\ref{cap:exponencials}: les construccions i els isomorfismes canònics que transporten les lleis.
\apartat{Objectius} En acabar el capítol, el lector hauria de poder:
\begin{itemize}
  \item tractar el cardinal \(\abs{A}\) com a classe d'equipotència, sense representant, i definir-ne la suma, el producte i l'exponenciació comprovant-ne la invariància;
  \item demostrar les lleis de l'aritmètica cardinal com a isomorfismes canònics, i la monotonia en l'exponent només per a base no nul·la;
  \item demostrar Cantor en forma cardinal i distingir \(2^{\aleph_0}\) de \(2^{\omega}\);
  \item demostrar que la comparabilitat total dels cardinals equival a l'elecció;
  \item construir el cardinal inicial d'un conjunt dins \(\Hart(A)\) i explicar la prudència amb \enquote{el cardinal és un conjunt};
  \item demostrar el teorema de König, identificant cada ús de l'elecció, i deduir-ne Cantor.
\end{itemize}
\end{objectius}

L'aritmètica ordinal recordava l'ordre; la cardinal l'oblida. Un
\emph{cardinal} és el que queda d'un conjunt quan se n'oblida qualsevol
ordre i es reté només la \emph{mida} ---la classe d'equipotència. Aquest
capítol desenvolupa l'aritmètica de les mides amb la mateixa disciplina de
sempre: les operacions \(+\), \(\cdot\) i l'exponenciació es defineixen
\emph{sobre representants} ---suma disjunta, producte i exponencial de
conjunts--- i es comprova que no depenen del representant, és a dir que estan
ben definides sobre classes d'equipotència. En aquest primer nivell no cal
cap \enquote{nombre cardinal} com a objecte: n'hi ha prou amb la relació
d'equipotència i la comparació per monomorfismes del capítol~\ref{cap:cantor}.
El teorema de Cantor hi reapareix com la desigualtat fonamental
\(\abs{A}<2^{\abs{A}}\). Després, amb l'elecció, veurem que dos conjunts
qualssevol són comparables, de manera que els cardinals queden totalment
ordenats; i que cada conjunt té un \emph{cardinal inicial} ---el menor tipus
d'ordre de la seva mida---, amb la prudència que aquest és un tipus, no un
conjunt escollit. Tanquem amb el teorema de König, un refinament diagonal de
Cantor per a famílies.

\section{Cardinalitat sense nombres cardinals}\label{sec:cardinalitat-relacional}

Recordem del capítol~\ref{cap:cantor} la comparació de conjunts: \(A\preceq B\)
si hi ha un monomorfisme \(A\hookrightarrow B\), \(A\approx B\) si hi ha un
isomorfisme (equipotència, definició~\ref{def:equipotencia}), i \(A\prec B\)
si \(A\preceq B\) però \(A\not\approx B\) (definició~\ref{def:comparacio}). El
\emph{cardinal} d'\(A\), que escrivim \(\abs{A}\), és \(A\) considerat llevat
d'equipotència; \(\abs{A}=\abs{B}\) vol dir \(A\approx B\). Com amb els tipus
ordinals, no cal escollir cap representant: treballem amb conjunts qualssevol
i comprovem la invariància.

\begin{definicio}[Operacions cardinals]\index{Operacions cardinals}\label{def:operacions-cardinals}
Per a conjunts \(A\) i \(B\), es defineixen la \emph{suma}, el \emph{producte}
i l'\emph{exponenciació} cardinals per
\[
\abs{A}+\abs{B}:=\abs{A+B},
\qquad
\abs{A}\cdot\abs{B}:=\abs{A\times B},
\qquad
\abs{B}^{\abs{A}}:=\abs{B^{A}},
\]
és a dir, com els cardinals de la suma disjunta
(definició~\ref{def:suma-disjunta}), del producte
(definició~\ref{def:producte}) i de l'exponencial
(definició~\ref{def:exponencial}), respectivament. Aquestes són operacions
\emph{sobre classes d'equipotència}, no sobre cap conjunt global de cardinals.
\end{definicio}

\begin{proposicio}[Les operacions estan ben definides]\label{prop:operacions-invariants}
La suma, el producte i l'exponenciació cardinals no depenen dels
representants: si \(A\approx A'\) i \(B\approx B'\), aleshores
\[
A+B\approx A'+B',
\qquad
A\times B\approx A'\times B',
\qquad
B^{A}\approx B'^{A'} .
\]
\emph{Axiomes usats:} suma, producte, exponencials
(capítols~\ref{cap:valors},~\ref{cap:parts} i~\ref{cap:exponencials}); cap
elecció.
\end{proposicio}

\begin{prova}
Siguin \(f\colon A\cong A'\) i \(g\colon B\cong B'\) isomorfismes. Aleshores
\(f+g\colon A+B\cong A'+B'\) (functorialitat de la suma disjunta,
capítol~\ref{cap:valors}) i \(f\times g\colon A\times B\cong A'\times B'\)
(functorialitat del producte, capítol~\ref{cap:parts}) són isomorfismes. Per a
l'exponencial, \(h\mapsto g\circ h\circ f^{-1}\) és una bijecció
\(B^{A}\cong B'^{A'}\), amb inversa \(k\mapsto g^{-1}\circ k\circ f\);
internament, és la transposada de
\(g\circ\operatorname{ev}\circ(\id_{B^A}\times f^{-1})\colon B^A\times A'\to B'\)
(propietat universal de l'exponencial, definició~\ref{def:exponencial}). Les
tres construccions són isomorfismes de conjunts, de manera que les mides
coincideixen.
\end{prova}

\begin{proposicio}[Lleis de l'aritmètica cardinal]\label{prop:lleis-cardinals}
Per a cardinals qualssevol \(\kappa=\abs{A}\), \(\mu=\abs{B}\),
\(\nu=\abs{C}\):
\begin{enumerate}
\item suma i producte són \emph{commutatius} i \emph{associatius}, amb neutres
\(\abs{\mathbf0}=0\) per a la suma i \(\abs{\mathbf1}=1\) per al producte;
\item distributivitat: \(\kappa\cdot(\mu+\nu)=\kappa\cdot\mu+\kappa\cdot\nu\);
\item lleis d'exponents: \(\kappa^{\mu+\nu}=\kappa^{\mu}\cdot\kappa^{\nu}\),
\((\kappa\cdot\mu)^{\nu}=\kappa^{\nu}\cdot\mu^{\nu}\),
\((\kappa^{\mu})^{\nu}=\kappa^{\mu\cdot\nu}\);
\item monotonia: si \(\kappa\le\mu\) (és a dir \(A\preceq B\)), aleshores
\(\kappa+\nu\le\mu+\nu\), \(\kappa\cdot\nu\le\mu\cdot\nu\) i
\(\kappa^{\nu}\le\mu^{\nu}\) (monotonia en la \emph{base}); a l'\emph{exponent},
si a més \(\nu\neq 0\), aleshores \(\nu^{\kappa}\le\nu^{\mu}\);
\item compatibilitat amb els recomptes finits: sobre conjunts finits, aquestes
operacions coincideixen amb la suma, el producte i la potència de naturals
(capítol~\ref{cap:finits}).
\end{enumerate}
\emph{Axiomes usats:} propietats universals de suma, producte i exponencial;
cap elecció.
\end{proposicio}

\begin{prova}
Cada llei és un isomorfisme canònic de conjunts, transportat a mides per la
proposició~\ref{prop:operacions-invariants}. \emph{(1)} La commutativitat de
la suma és \(A+B\cong B+A\) i la del producte \(A\times B\cong B\times A\); les
associativitats són els isomorfismes canònics corresponents; els neutres són
\(A+\mathbf0\cong A\) i \(A\times\mathbf1\cong A\). \emph{(2)} És la
distributivitat \(A\times(B+C)\cong(A\times B)+(A\times C)\)
(capítol~\ref{cap:valors}). \emph{(3)} Són els isomorfismes exponencials, que alineem lletra a lletra amb
l'enunciat. Per a \(\kappa^{\mu+\nu}=\kappa^{\mu}\cdot\kappa^{\nu}\), amb
\(\kappa=\abs C\), \(\mu=\abs B\), \(\nu=\abs A\): \(C^{B+A}\cong C^{B}\times
C^{A}\) (una aplicació d'una suma és un parell d'aplicacions). Per a
\((\kappa\cdot\mu)^{\nu}=\kappa^{\nu}\cdot\mu^{\nu}\), amb \(\kappa=\abs C\),
\(\mu=\abs B\), \(\nu=\abs A\): \((C\times B)^{A}\cong C^{A}\times B^{A}\). Per
a \((\kappa^{\mu})^{\nu}=\kappa^{\mu\cdot\nu}\), amb \(\kappa=\abs C\),
\(\mu=\abs B\), \(\nu=\abs A\): \((C^{B})^{A}\cong C^{B\times A}\)
(currificació, definició~\ref{def:exponencial}). En els tres casos les lletres
dels representants s'assignen igual que a l'enunciat, i l'ordre del producte
\(B\times A\) és el que correspon a \(\mu\cdot\nu\) (indiferent per
commutativitat).
\emph{(4)} La monotonia \emph{en la base} surt de compondre un monomorfisme
\(A\hookrightarrow B\) amb les operacions, que preserven monomorfismes: de
\(A\hookrightarrow B\) s'obté \(A+C\hookrightarrow B+C\),
\(A\times C\hookrightarrow B\times C\) i \(A^{C}\hookrightarrow B^{C}\) per
postcomposició. Per a l'\emph{exponent}, en canvi, un monomorfisme
\(A\hookrightarrow B\) indueix per \emph{precomposició} una aplicació
\(C^{B}\to C^{A}\) ---en la direcció contrària---, i per anar de \(C^{A}\) cap
a \(C^{B}\) cal estendre cada \(f\colon A\to C\) a \(B\to C\). Amb
\(\nu=\abs C\neq 0\) fixem un punt \(c_0\in C\) i estenem per casos: sobre la
imatge de \(A\) copiem \(f\), i sobre el complement de \(A\) dins \(B\)
posem el valor constant \(c_0\); això dona un monomorfisme \(C^{A}\hookrightarrow
C^{B}\), és a dir \(\nu^{\kappa}\le\nu^{\mu}\). Quan \(\nu=0\) la construcció
falla i l'enunciat és fals: \(0^{0}=1\not\le 0=0^{1}\), tot i que \(0\le 1\).
Per això la monotonia en l'exponent només s'afirma per a \(\nu\neq 0\).
\emph{(5)} Sobre segments
\(\seg{m}\) i \(\seg{n}\), \(\seg{m}+\seg{n}\approx\seg{m+n}\),
\(\seg{m}\times\seg{n}\approx\seg{m\cdot n}\) i
\(\seg{n}^{\seg{m}}\approx\seg{n^{m}}\) són els recomptes finits del
capítol~\ref{cap:finits}.
\end{prova}

\begin{teorema}[Cantor, en forma cardinal]\label{teo:cantor-cardinal}
Per a tot conjunt \(A\),
\[
\abs{A}\ <\ 2^{\abs{A}} .
\]
És a dir, \(\abs{A}\le 2^{\abs{A}}\) i \(\abs{A}\neq 2^{\abs{A}}\), on
\(2^{\abs{A}}=\abs{\mathbf2^{A}}=\abs{P(A)}\).\par
\emph{Axiomes usats:} teorema de Cantor (teorema~\ref{teo:cantor}), exponencial
\(P(A)\cong\mathbf2^{A}\) (capítol~\ref{cap:exponencials}); cap elecció.
\end{teorema}

\begin{prova}
Com que \(P(A)\cong\mathbf2^{A}\) (el conjunt de parts és l'exponencial de
base \(\mathbf2\), capítol~\ref{cap:exponencials}), es té
\(2^{\abs{A}}=\abs{\mathbf2^{A}}=\abs{P(A)}\). El teorema de Cantor
(teorema~\ref{teo:cantor}) diu \(A\prec P(A)\), és a dir
\(\abs{A}\le\abs{P(A)}\) i \(\abs{A}\neq\abs{P(A)}\). Per tant
\(\abs{A}<2^{\abs{A}}\).
\end{prova}

\begin{observacio}[Ordinal contra cardinal, un cop més]\label{obs:cardinal-vs-ordinal}
L'exponenciació cardinal \(2^{\abs{A}}\) és radicalment diferent de l'ordinal
del capítol anterior. Ordinalment, \(2^{\omega}=\omega\)
(proposició~\ref{prop:exponenciacio-calculs}): una longitud numerable.
Cardinalment, \(2^{\aleph_0}=\abs{P(\N)}>\aleph_0\) pel teorema de Cantor: una
mida no numerable, el continu. La mateixa notació \enquote{\(2\) elevat a}
significa coses oposades segons si recorda l'ordre (ordinal) o l'oblida
(cardinal). És per evitar aquesta confusió que hem construït les dues
aritmètiques per separat, i que escriurem \(\aleph_0\) per al cardinal de
\(\N\) i \(\omega\) per al seu tipus d'ordre, tot i que \enquote{són} el
mateix conjunt.
\end{observacio}

\section{Comparabilitat i ordre total dels cardinals}\label{sec:comparabilitat-cardinals}

En el primer nivell, la comparació \(\preceq\) és un preordre invariant
(proposició~\ref{prop:preceq-preordre}) i, per Schröder--Bernstein, és
antisimètrica mòdul equipotència (corol·lari~\ref{cor:preceq-antisimetric}).
El que \emph{no} tenim sense elecció és que sigui \emph{total}: que dos
conjunts qualssevol siguin comparables. Això ho dona el teorema del bon ordre.

\begin{teorema}[Comparabilitat total dels cardinals]\label{teo:comparabilitat-cardinals}
Suposem l'axioma d'elecció (axioma~\ref{ax:eleccio}). Aleshores dos conjunts
qualssevol són comparables:
\[
A\preceq B
\qquad\text{o bé}\qquad
B\preceq A .
\]
Combinat amb Schröder--Bernstein (teorema~\ref{teo:schroder-bernstein}), la
comparació de cardinals és un \emph{ordre total llevat d'equipotència}:
reflexiu, transitiu, antisimètric mòdul \(\approx\) i total.\par
\emph{Axiomes usats:} teorema del bon ordre (teorema~\ref{teo:bon-ordre}) ---i
per tant elecció---, comparabilitat de bons ordres
(teorema~\ref{teo:comparabilitat}), Schröder--Bernstein
(teorema~\ref{teo:schroder-bernstein}).
\end{teorema}

\begin{prova}
Pel teorema del bon ordre (teorema~\ref{teo:bon-ordre}), \(A\) i \(B\) admeten
bons ordres; siguin \(W_A\) i \(W_B\) representants. Per la comparabilitat de
bons ordres (teorema~\ref{teo:comparabilitat}), un dels dos és isomorf a un
segment inicial de l'altre ---diguem \(W_A\) isomorf a un segment de \(W_B\).
Aquest isomorfisme d'ordre és en particular un monomorfisme de conjunts
\(A\hookrightarrow B\), d'on \(A\preceq B\). Simètricament en l'altre cas. Per
tant \(A\preceq B\) o \(B\preceq A\). La resta ---reflexivitat, transitivitat,
antisimetria mòdul \(\approx\)--- és la proposició~\ref{prop:preceq-preordre} i
el corol·lari~\ref{cor:preceq-antisimetric}.
\end{prova}

\begin{observacio}[La comparabilitat és equivalent a l'elecció]\label{obs:comparabilitat-eleccio}
La comparabilitat total dels cardinals no és només \emph{conseqüència} de
l'elecció: n'és \emph{equivalent}. Si dos conjunts qualssevol són comparables,
aleshores en particular tot conjunt \(A\) és comparable amb \(\Hart(A)\); com
que \(\Hart(A)\not\preceq A\) (teorema~\ref{teo:hartogs}), ha de ser
\(A\preceq\Hart(A)\), és a dir \(A\) s'injecta en un bon ordre, i per tant
hereta un bon ordre ---d'on el teorema del bon ordre, equivalent a l'elecció
(teorema~\ref{teo:equivalencies}). Aquesta és una altra manifestació del
paper de Hartogs: la seva cota converteix la comparabilitat en bon ordre.
\end{observacio}

\section{Cardinals inicials}\label{sec:cardinals-inicials}

Fins ara \enquote{\(\abs{A}\)} ha estat una classe d'equipotència, no un
objecte. Amb l'elecció podem, si volem, triar per a cada mida un
\emph{representant} distingit: el menor tipus d'ordre d'aquella mida. Cal
fer-ho amb la prudència que la guia recomana.

\begin{definicio}[Bon ordre inicial i cardinal inicial]\index{Bon ordre inicial i cardinal inicial}\label{def:cardinal-inicial}
Un bon ordre \(W\) és \emph{inicial} si cap bon ordre de tipus estrictament
menor no és equipotent a \(W\): per a tot bon ordre \(V\) amb
\(\operatorname{tipus}(V)<\operatorname{tipus}(W)\), es té \(V\not\approx W\).
Amb l'elecció, cada conjunt \(A\) és equipotent a algun bon ordre
(teorema~\ref{teo:bon-ordre}); el \emph{cardinal inicial} de \(A\) és el
\emph{menor} tipus d'ordre \(\alpha\) tal que un (tot) representant seu és
equipotent a \(A\). Aquest tipus és inicial i únic.
\end{definicio}

\begin{proposicio}[Existència i unicitat del cardinal inicial]\label{prop:cardinal-inicial}
Suposem l'elecció. Per a tot conjunt \(A\) existeix un únic cardinal inicial
\(\alpha\), i és un tipus d'ordre inicial. La cerca està \emph{acotada}: el
menor tipus equipotent a \(A\) és un element de \(\Hart(A)\), sense recórrer
cap classe global.\par
\emph{Axiomes usats:} teorema del bon ordre (teorema~\ref{teo:bon-ordre}),
Hartogs i el bon ordre de \(\Hart\) (proposició~\ref{prop:hartogs-bon-ordre});
elecció.
\end{proposicio}

\begin{prova}
Sigui \(E\hookrightarrow\Hart(A)\) la part dels tipus \([(S,R)]\) amb
\(S\approx A\) (una condició interna sobre \(\WO(A)\), invariant per
isomorfisme, que baixa al quocient). És no inicial: pel teorema del bon ordre
\(A\) admet un bon ordre \(R\), i \([(A,R)]\in E\). Per ser \(\Hart(A)\) un bon
ordre (proposició~\ref{prop:hartogs-bon-ordre}), \(E\) té un mínim \(\alpha\),
únic. És inicial: si un bon ordre \(V\) amb \(\operatorname{tipus}(V)<\alpha\)
fos equipotent a \(A\), aleshores \(V\) seria isomorf a un segment propi d'un
representant \((S,R)\) de \(\alpha\), és a dir, a un bon ordre sobre una part
de \(S\subseteq A\); el seu tipus seria un element de \(E\) menor que
\(\alpha\), contra la minimalitat. Tota la cerca ha passat dins \(\Hart(A)\),
acotadament.
\end{prova}

\begin{observacio}[\enquote{El cardinal és un conjunt}: una prudència]\label{obs:cardinal-prudencia}
Convé ser prudent amb la frase \enquote{el cardinal de \(A\) és un conjunt}.
En el tractament estructural, el cardinal inicial de \(A\) és un \emph{tipus
d'ordre} inicial ---únic llevat d'un únic isomorfisme d'ordre
(teorema~\ref{teo:unicitat-iso})---, no un conjunt privilegiat escollit
d'entrada. Triar-ne un representant concret (un bon ordre determinat d'aquell
tipus) és una \emph{tria addicional}, legítima però no canònica. Aquesta és la
raó per la qual, en el primer nivell, hem treballat amb \(\abs{A}\) com a
classe d'equipotència i no com a objecte: la teoria relacional no necessita
representants, i només els introdueix ---amb l'elecció--- qui vulgui una
successió distingida de cardinals. Una \emph{successió global} de tots els
cardinals inicials, d'altra banda, seria de Nivell~II: reunir-los tots
demanaria SR (observació~\ref{obs:sr}), igual que la successió dels
\(\aleph_\alpha\) que veurem al capítol opcional.
\end{observacio}

\section{El teorema de König}\label{sec:konig}

Tanquem amb un refinament de Cantor per a famílies: si cada terme d'una
família és estrictament més petit que el corresponent d'una altra, la suma
dels primers és estrictament menor que el producte dels segons. La prova és,
en el fons, la diagonal de Cantor.

\begin{teorema}[König]\index{König, teorema de}\label{teo:konig}
Suposem l'elecció. Siguin \((A_i)_{i\in I}\) i \((B_i)_{i\in I}\) dues famílies
de conjunts, indexades pel mateix \(I\), amb \(A_i\prec B_i\) per a cada \(i\).
Aleshores
\[
\sum_{i\in I}\abs{A_i}\ <\ \prod_{i\in I}\abs{B_i},
\]
on la suma és el cardinal de la suma disjunta \(\coprod_i A_i\) i el producte
el del producte \(\prod_i B_i\) (capítol~\ref{cap:families}).\par
\emph{Axiomes usats:} elecció (axioma~\ref{ax:eleccio}), famílies, suma i
producte (capítol~\ref{cap:families}), la diagonal de Cantor
(teorema~\ref{teo:diagonal}).
\end{teorema}

\begin{prova}
L'elecció s'usa en tres llocs, que assenyalem: per triar alhora un
monomorfisme \(m_i\) i un punt fora de la seva imatge per a cada \(i\); per
deduir de \(F_i\) exhaustiva que \(B_i\preceq A_i\) (una secció); i per triar
alhora els punts \(d_i\) de la diagonal. Les tres tries són seccions de
famílies internes indexades per \(I\), construïdes amb exponencials
(capítol~\ref{cap:families}).

\emph{\(\le\).} Hi ha un monomorfisme \(\coprod_i A_i\hookrightarrow\prod_i B_i\):
com que \(A_i\prec B_i\), en particular \(A_i\preceq B_i\), i triant (per
l'elecció) un monomorfisme \(m_i\colon A_i\hookrightarrow B_i\) i un punt
\(b_i\in B_i\setminus\operatorname{Im}(m_i)\) ---no buit perquè \(m_i\) no és
exhaustiu, ja que \(A_i\not\approx B_i\)---, enviem un element \(a\) de la
fibra \(A_i\) a la família \((x_j)_{j\in I}\) amb \(x_i=m_i(a)\) i \(x_j=b_j\)
per a \(j\neq i\). És injectiva perquè la coordenada \(i\) i la posició del
\enquote{no-\(b\)} determinen \(a\) i \(i\).

\emph{\(\neq\).} Vegem que cap aplicació \(F\colon\coprod_i A_i\to\prod_i B_i\)
no és exhaustiva ---la diagonal. Per a cada \(i\), la restricció de \(F\) a la
fibra \(A_i\), seguida de la projecció \(\pi_i\colon\prod_j B_j\to B_i\), dona
una aplicació \(F_i\colon A_i\to B_i\); com que \(A_i\prec B_i\), \(F_i\) no és
exhaustiva (si ho fos, una secció seva, que existeix per l'elecció, donaria
\(B_i\preceq A_i\), i per Schröder--Bernstein \(A_i\approx B_i\)), de manera que existeix
\(d_i\in B_i\setminus\operatorname{Im}(F_i)\) ---triem-ne un per l'elecció. La
família \((d_i)_{i\in I}\in\prod_i B_i\) no és cap valor de \(F\): si fos
\(F(a)=(d_i)_i\) amb \(a\) a la fibra \(A_{i_0}\), aleshores la coordenada
\(i_0\) donaria \(F_{i_0}(a)=d_{i_0}\), contra \(d_{i_0}\notin\operatorname{Im}(F_{i_0})\).
Per tant \(F\) no és exhaustiva, cap \(\coprod_i A_i\to\prod_i B_i\) no ho és,
i \(\sum_i\abs{A_i}\neq\prod_i\abs{B_i}\). Amb la desigualtat \(\le\), s'obté
l'estricta.
\end{prova}

\begin{observacio}[König conté Cantor]\label{obs:konig-cantor}
El teorema de Cantor és el cas particular de König amb \(I=A\), \(A_i=\mathbf1\)
i \(B_i=\mathbf2\) per a tot \(i\): aleshores \(A_i=\mathbf1\prec\mathbf2=B_i\),
i König dona
\[
\abs{A}=\sum_{i\in A}1<\prod_{i\in A}2=2^{\abs{A}},
\]
que és exactament \(\abs{A}<2^{\abs{A}}\) (teorema~\ref{teo:cantor-cardinal}).
La \enquote{idea diagonal} que apareix a la prova de König ---construir una
família \((d_i)\) que difereix de cada valor en alguna coordenada--- és
literalment la diagonal de Cantor, ara repartida coordenada a coordenada. Per
això König es pot veure com la forma \emph{indexada} del teorema de Cantor.
\end{observacio}

\begin{resumcap}
\apartat{Idees centrals}
\begin{enumerate}
  \item El cardinal \(\abs{A}\) és la classe del conjunt \(A\) llevat d'equipotència (no un tipus d'ordre: no hi ha ordre); suma, producte i exponenciació són els cardinals de \(A+B\), \(A\times B\) i \(B^A\), i les lleis són isomorfismes canònics. La monotonia en l'exponent demana base no nul·la: \(0^0=1\not\le0=0^1\).
  \item Cantor: \(\abs{A}<2^{\abs{A}}\). L'exponenciació cardinal \(2^{\aleph_0}\) no té res a veure amb l'ordinal \(2^\omega=\omega\).
  \item Sense elecció, \(\preceq\) és un preordre antisimètric mòdul \(\approx\); la totalitat és equivalent a l'elecció (Hartogs fa el pas de comparabilitat a bon ordre).
  \item Amb l'elecció, el cardinal inicial d'\(A\) és el menor tipus d'ordre de \(\Hart(A)\) el suport d'un representant del qual és equipotent a \(A\); és un tipus, no un conjunt, i triar-ne un representant és una tria més.
  \item König: si \(A_i\prec B_i\) per a tot \(i\), \(\sum\abs{A_i}<\prod\abs{B_i}\). La prova és una diagonal coordenada a coordenada, i usa l'elecció tres vegades. Conté Cantor com a cas particular.
\end{enumerate}
\apartat{Errors típics} Escriure \(\abs{A}\) com si fos un conjunt fix; aplicar la monotonia en l'exponent amb base \(0\); confondre \(2^\omega\) ordinal amb \(2^{\aleph_0}\) cardinal; creure que la comparabilitat de cardinals és gratuïta; i passar de \enquote{exhaustiva} a \enquote{domina} sense elecció.
\apartat{Lectura recomanada} \cite[§§22--25]{halmos} per a l'aritmètica cardinal clàssica; \cite[cap.~3 i~5]{jech2003} per als cardinals, König i el paper de l'elecció; \cite[cap.~I.10]{kunen2011} per als cardinals inicials; i \cite[§§8.5--8.6]{lawvere-rosebrugh} per a la mida dels conjunts abstractes.
\end{resumcap}
